\paragraph{Bounds without the learned prior.} A reconstruction selects one volume, even when several volumes are consistent with the measurements. To assess this ambiguity without relying on the learned reconstruction prior, we use set-membership estimation \cite{milanese1991,haldar2023}. Given the intensity bounds $\bx\in[0,1]^N$, we define the measurement-consistency set
\begin{equation*}
\C(\by)=\{\bx\in[0,1]^N:\ \ell_i\le(\bA\bx)_i\le h_i\},
\end{equation*}
where the intervals $[\ell_i,h_i]$ jointly contain the noise-free measurements with a stated probability under the assumed noise model. For a non-empty consistency set, minimising and maximising $x_j$ over $\C(\by)$ gives the smallest and largest intensities consistent with these constraints. Both bounds are computed by linear programming. Their difference measures the remaining ambiguity in that voxel's intensity.

The connection to optical learning can be understood through a two-voxel example. Let $\mathbf a_1,\mathbf a_2$ be unit-norm measurement columns with correlation $\rho$, where $|\rho|<1$. For coefficient perturbations satisfying $\|\mathbf a_1\delta x_1+\mathbf a_2\delta x_2\|_2\le\varepsilon$, allowing $\delta x_2$ to vary freely gives $|\delta x_1|\le\varepsilon/\sqrt{1-\rho^2}$. Highly correlated responses therefore permit greater ambiguity. This provides intuition for reducing depth correlations, rather than a guarantee for the full system, which couples many voxels, includes intensity constraints, and uses raw measurement columns rather than the centred training signatures. We compute the actual bounds on reduced fields where the linear programs are tractable.
